\documentclass[11pt]{article}

\usepackage[margin=1in]{geometry}
\usepackage{amsmath,amssymb,amsthm}
\usepackage{booktabs,array}
\usepackage{xcolor}
\usepackage[hidelinks]{hyperref}

\newtheorem{theorem}{Theorem}[section]
\newtheorem{proposition}[theorem]{Proposition}
\newtheorem{lemma}[theorem]{Lemma}
\newtheorem{corollary}[theorem]{Corollary}
\newtheorem{conjecture}[theorem]{Conjecture}
\theoremstyle{definition}
\newtheorem{definition}[theorem]{Definition}
\theoremstyle{remark}
\newtheorem{remark}[theorem]{Remark}

\newcommand{\F}{\mathbb F}
\newcommand{\supp}{\operatorname{supp}}
\newcommand{\bd}{\partial}
\newcommand{\Good}{\mathcal G}

\title{Two connected countermodels to one-circuit repair\\
for five-cycle double covers}
\author{AI-assisted research draft for independent human verification\\
\small Human author attribution to be supplied before submission}
\date{26 July 2026}

\begin{document}
\maketitle

\noindent
\colorbox{yellow!22}{%
  \parbox{\dimexpr\linewidth-2\fboxsep\relax}{%
  \textbf{AI-assistance disclosure.}
  OpenAI Codex agents (GPT-5-series models, accessed 25--26 July 2026)
  proposed the successful-value and one-circuit-repair formulations,
  formulated the pure-coordinate-merge strategy,
  generated proof ideas and substantial prose, wrote the producer and
  checking programs, performed the computational searches and literature
  audit, found the finite countermodels, designed the two-sum
  constructions, discovered the 40- and 46-vertex reductions and the
  rigid-cap obstruction, and generated the checking artifacts, all under
  human direction.  A second Codex agent
  was asked to audit the logic
  adversarially.  The ``independent'' Python checker, the separately
  written C++ implementation, and that audit were also AI-assisted; none
  is independent human verification or peer review.  The main results
  refute two intermediate one-circuit strategies, not the five-cycle
  double cover conjecture; both constructed graphs positively have
  standard five-cycle double covers.  No human author or referee is represented as
  having certified the manuscript, proof, certificate, or computation.
  This research draft is not journal-ready and should not be submitted
  until human authors check every proof and incidence table, independently
  reconstruct the graph, rerun the computations, verify the literature
  attribution, and accept ordinary scholarly responsibility.}}

\begin{abstract}
We give two explicit countermodels to one-circuit repair strategies for
the five-cycle double cover conjecture.  The first uses a
nowhere-zero \(\F_2^3\)-flow \(f\) on a simple connected bridgeless cubic
40-vertex graph.  No value class of \(f\), or of any allowed neighbour
obtained by switching one connected circuit, satisfies a particular
edge-disjoint \(T\)-join packing condition sufficient for a 5-CDC.

The second concerns the potential construction of an eight-cover from
a fixed \(\F_2^3\)-flow used in recent proofs and expositions of the
cycle double cover theorem.  On an explicit simple connected bridgeless
cubic 46-vertex graph, no compatible eight-cover for the displayed flow
can be purely merged into at most five Eulerian subgraphs.  The same
remains true after every nowhere-zero switch on one connected circuit.
A 12-vertex port-deleted block is potential-rigid up to coordinate
translation and forces a \(K_6\) in the coordinate co-occurrence graph.
Three copies are attached at three same-colour base edges that no circuit
can contain simultaneously; a short triangle argument proves this
non-cyclability.  These ingredients give SAT-free proofs of both global
one-switch assertions.

Both countermodels are Tait-colourable and have explicit three-cycle
double covers.  They therefore refute only the stated intermediate
strategies and do not resolve or refute the five-cycle double cover
conjecture.

For context, exact enumeration had found no obstruction through all
connected simple bridgeless cubic graphs of order at most 14, selected
complete hard-graph scopes through order 20, and all 20 strict snarks of
order 22 and the 38 strict order-24 snarks output by a retained
Snarkhunter 2.0b run.  Connected-circuit enumeration checks
\(847{,}539{,}220\) flow-orbit representatives at order 20,
\(63{,}251{,}330\) at order 22, and \(617{,}079{,}260\) at order 24.
The 40-vertex construction shows that these are finite phenomena, not a
universal theorem.  All artifact generation, checking code, and proof
development were AI-assisted and still require independent human
verification.
\end{abstract}

\section{Introduction and attribution}

A \emph{standard five-cycle double cover} (5-CDC) of a finite graph is a
list \(C_1,\ldots,C_5\) of Eulerian edge-subsets such that every edge
belongs to exactly two members.  Here Eulerian means even degree at every
vertex; connectedness is not required.  Empty members are permitted, so a
cover with at most five members can be padded to five.  This note concerns
the standard, unoriented conjecture only.  This agrees with the current
``at most \(k\)'' convention in Oum's exposition
\cite[Section~9.1]{oum2026}.

The existence of a 5-CDC for every bridgeless graph remains unproved.
Oum's July 2026 exposition gives an 8-CDC and records the 5-CDC statement
separately as Conjecture~18 \cite[Section~9.1]{oum2026}.  See also the
contemporary expositions of Geelen and Kintali
\cite{geelen2026,kintali2026}.  The
present note does not claim otherwise.  Its two main results instead
refute two specific intermediate proof strategies inside the space of
nowhere-zero \(\F_2^3\)-flows.  Both countermodels have explicitly
checked 5-CDCs.

The underlying matching/nowhere-zero-4-flow characterization is prior:
Hoffmann-Ostenhof proves that a cubic graph has a 5-CDC precisely when it
has a matching \(M\), two Eulerian edge-subsets intersecting exactly in
\(M\), and a nowhere-zero 4-flow on \(G-M\)
\cite[Corollary~0.6]{hoffmann2013}.  We restate and prove the exact
specialization needed below to eliminate convention ambiguity.
Further 4-flow approaches and sufficient conditions for 5-CDCs, including
Catlin-reduction and weak-oddness results, are due to Liu, Hao, Luo, and
Zhang \cite{liu2023,liu2025}.  None of those results is claimed as new
here.

The reconfiguration language is also prior.  Esperet, Hendrey, Lagoutte,
Marseloo, Norin, and Steiner define the graph of nowhere-zero flows in
which consecutive flows differ on one circuit and prove positive and
negative connectivity results \cite{esperet2025}.  Cranston, Li, Su,
Wang, and Xu develop this theory further, including reductions to cubic
graphs and results for several groups and graph classes
\cite{cranston2026}.  In particular, the definition of circuit-supported
adjacency is not a contribution of this note.  More precisely, Esperet
et al.\ prove that for every 2-edge-connected graph the reconfiguration
graph has minimum degree at least one when the abelian flow group has
order at least six or is \(\mathbb Z_2^2\); they also treat
\(\mathbb Z_4\), while frozen \(\mathbb Z_5\)-flows exist
\cite{esperet2025}.  Even the positive minimum-degree statements are much
weaker than requiring every unsuccessful flow to have a
\emph{successful} neighbour.  Their universal vector-space connectivity
result uses \(\mathbb Z_2^8\), not \(\F_2^3\).

The term \emph{Fano-flow conjectures} also has an established, different
meaning: Jin, Mazzuoccolo, and Steffen study conjectures controlling how
many lines of the Fano plane are used by a flow
\cite{jin2018}.  Our successful-value condition is instead a
value-class/\(T\)-join packing condition, and \((\mathsf D)\) is a
domination statement in a reconfiguration graph.  Neither countermodel
below refutes any of those classical Fano-flow conjectures.

We found no statement in the sources above of either specific
one-circuit domination assertion, the pure-merge coloring obstruction,
or the two countermodels proved here.  Accordingly, the only priority
claim made here is provisional: these two negative results appear not to
have been recorded.  The successful-value definition, exact blocks,
two-sum compositions, tetrahedral \(T\)-join observation, and finite
computations belong to those constructions; the underlying flow,
\(T\)-join, two-sum, potential, and circuit-reconfiguration machinery is
not claimed as new.  In particular, the flow-to-pair-label potential
construction appears in Oum's and Kintali's presentations
\cite[Section~5]{oum2026} and \cite{kintali2026}.  An independent expert literature
review, including consultation with authors of the recent flow and
reconfiguration papers, is required before submission.

\section{Pair labels and the exact Boolean encoding}
\label{sec:labels}

We first fix the direct 5-CDC convention independently of the later flow
reformulation.  For this section an undirected loop contributes its two
incidences at its vertex.  The remaining sections and all computations
are loopless.

\begin{proposition}[pair-label equivalence]\label{prop:labels}
For a finite undirected graph \(G\), the following are equivalent.
\begin{enumerate}
  \item \(G\) has Eulerian edge-subsets \(C_1,\ldots,C_5\) and every edge
        belongs to exactly two of them.
  \item Every edge \(e\) has a label
        \(L(e)\in\binom{[5]}2\), and for every vertex \(v\) and
        coordinate \(i\in[5]\), coordinate \(i\) occurs evenly among the
        labels on the edge incidences at \(v\).
\end{enumerate}
\end{proposition}

\begin{proof}
Given the cover, put \(L(e)=\{i:e\in C_i\}\).  Exact double coverage gives
\(|L(e)|=2\), and the Eulerian condition on \(C_i\) gives the required
incidence parity.  Conversely, put \(C_i=\{e:i\in L(e)\}\).  The incidence
parity makes every \(C_i\) Eulerian, and \(|L(e)|=2\) makes every edge
belong to exactly two members.
\end{proof}

This equivalence has a literal SAT/XOR encoding.  Introduce
\(x_{e,i}\in\{0,1\}\), with \(x_{e,i}=1\) meaning \(i\in L(e)\), and impose
\[
 \sum_{i=1}^5x_{e,i}=2\quad(e\in E(G)),\qquad
 \bigoplus_{e\ni v}x_{e,i}=0\quad(v\in V(G),\ i\in[5]).             \tag{1}
\]
Incidences, rather than merely edge names, are counted in the second
equation.  Proposition~\ref{prop:labels} proves that (1) is equivalent to
the intended definition.

For completeness, the exactly-two condition on five variables has the
certificate-producing CNF
\[
 \neg x_i\vee\neg x_j\vee\neg x_k
 \quad(1\le i<j<k\le5),                            \tag{2}
\]
together with
\[
 \bigvee_{j\ne i}x_j\quad(i\in[5]).                \tag{3}
\]
The ten clauses (2) give at most two true variables.  If zero variables
are true, every clause (3) fails; if only \(x_i\) is true, the clause
omitting \(i\) fails.  Hence (2)--(3) give exactly two.

At a cubic vertex, writing one coordinate on its incident edges as
\(a,b,c\), even parity is exactly the four-clause CNF
\[
 (a\vee b\vee\neg c),\quad(a\vee\neg b\vee c),\quad
 (\neg a\vee b\vee c),\quad(\neg a\vee\neg b\vee\neg c).           \tag{4}
\]
These clauses exclude the four odd assignments.  Equations (1), or their
CNF expansion (2)--(4), are the authoritative standard-cover semantics;
no orientation variables occur.

\section{Value classes and \texorpdfstring{\(T\)}{T}-joins}
\label{sec:value}

From now on \(G\) is a finite loopless cubic graph and
\[
 f:E(G)\longrightarrow A-\{0\},\qquad A=\F_2^3,
\]
is a nowhere-zero flow.  Since \(A\) has exponent two, edge orientations
play no role: at every vertex the xor of the three incident values is
zero.  For \(k\in A-\{0\}\), put
\[
 M_k=\{e:f(e)=k\}.
\]
For an edge set \(X\), write \(\bd X\) for the set of vertices incident
with an odd number of edges of \(X\).  If \(T\subseteq V(H)\), a
\emph{\(T\)-join} of \(H\) is an edge set \(J\) with \(\bd J=T\).

\begin{proposition}[value-class structure]\label{prop:value}
For every \(k\ne0\):
\begin{enumerate}
  \item \(M_k\) is a matching;
  \item the quotient of \(f\) by \(\langle k\rangle\) is an
        \(\F_2^2\)-flow whose exact zero set is \(M_k\);
  \item \(G-M_k\) has a canonical family of four
        \(\bd M_k\)-joins, indexed by the linear functionals
        \[
          \Lambda_k=\{\lambda\in A^*:\lambda(k)=1\};
        \]
  \item these four joins double-cover \(G-M_k\), and each unordered pair
        of joins intersects in one of the other six value classes.
\end{enumerate}
\end{proposition}

\begin{proof}
At a cubic vertex the three incident nonzero values sum to zero.  Two
cannot be equal, since the third would then be zero.  Thus a fixed value
occurs at most once at a vertex, proving that \(M_k\) is a matching.

The zero coset of \(A/\langle k\rangle\) is \(\{0,k\}\).  Since \(f\) is
nowhere zero, an edge maps to zero in the quotient exactly when its
original value is \(k\).  This proves the second assertion.

There are four linear functionals in \(\Lambda_k\).  For
\(\lambda\in\Lambda_k\), the support of the binary flow
\(\lambda\circ f\) is an Eulerian edge set containing \(M_k\).  Therefore
\[
 J_\lambda=\{e\notin M_k:\lambda(f(e))=1\}
\]
has boundary \(\bd M_k\) in \(G-M_k\).

If \(f(e)=w\notin\{0,k\}\), the two independent equations
\(\lambda(k)=\lambda(w)=1\) have exactly two solutions in the
three-dimensional dual space.  Hence \(e\) occurs in exactly two of the
four joins.  Finally, if \(\lambda\ne\mu\) belong to \(\Lambda_k\), the
simultaneous equations
\[
 \lambda(x)=\mu(x)=1
\]
have exactly two solutions in \(A\): \(k\) and one value
\(w\notin\{0,k\}\).  After \(M_k\) is deleted,
\(J_\lambda\cap J_\mu=M_w\).  The six unordered pairs of functionals give
the six values \(w\ne0,k\).
\end{proof}

In particular, the elementary component-parity obstruction to the
existence of even one \(T\)-join never occurs for a value class.  The
issue is whether two such joins can be made edge-disjoint; the four
canonical joins need not contain a disjoint pair.

\begin{lemma}[even-marked circuit criterion]\label{lem:criterion}
Let \(M\) be a matching in a cubic graph \(G\), put
\(K=G-M\) and \(T=\bd M\).  Then \(K\) has two edge-disjoint
\(T\)-joins if and only if \(K\) has an Eulerian edge set \(Q\) that
contains every vertex of \(T\) and whose every nonempty circuit component
contains an even number of vertices of \(T\).
\end{lemma}

\begin{proof}
Suppose \(J_1,J_2\) are edge-disjoint \(T\)-joins.  A terminal has degree
two in \(K\); each join has odd degree there, so each uses one edge, and
disjointness makes them use the two different edges.  At a nonterminal,
each join has even degree in a three-edge incidence set.  Thus
\(Q=J_1\mathbin{\dot\cup}J_2\) has degree two at every terminal and degree
zero or two elsewhere.  It is a disjoint union of circuits.  Along a
circuit, the join colour stays unchanged at a nonterminal and switches at
a terminal.  The colour can close consistently only after an even number
of switches, so every circuit contains evenly many terminals.

Conversely, traverse each circuit of \(Q\), colouring its edges with two
join colours.  Keep the colour through a nonterminal and switch it at a
terminal.  The even terminal count makes the colouring consistent when
the traversal closes.  Each colour class has degree one at every terminal
and degree zero or two elsewhere, so the two colour classes are
edge-disjoint \(T\)-joins.
\end{proof}

\begin{definition}
A value \(k\ne0\) is \emph{successful} for \(f\) if
\((G-M_k,\bd M_k)\) has two edge-disjoint \(T\)-joins.  The flow \(f\) is
\emph{good} if it has at least one successful value.
\end{definition}

\begin{proposition}[successful value implies a 5-CDC]\label{prop:success}
If \(f\) is good, then \(G\) has a standard five-cycle double cover.
\end{proposition}

\begin{proof}
Let \(k\) be successful, put \(M=M_k\), and take edge-disjoint
\(T\)-joins \(J_1,J_2\) in \(G-M\), where \(T=\bd M\).  Then
\[
 A_1=M\mathbin{\dot\cup}J_1,\qquad
 A_2=M\mathbin{\dot\cup}J_2
\]
are Eulerian edge sets of \(G\), and \(A_1\cap A_2=M\).

Let \(\phi:E(G)\to\F_2^2\) be the quotient flow from
Proposition~\ref{prop:value}; its exact zero set is \(M\).  On
\(K=G-M\), put
\[
 D_0=A_1\mathbin{\triangle}A_2=J_1\mathbin{\dot\cup}J_2.
\]
We construct a four-cycle double cover of \(K\) whose first coordinate is
\(D_0\).  Represent the nonzero elements \(1,2,3\) of \(\F_2^2\) by the
linear section
\[
 s(1)=1100,\qquad s(2)=1010,\qquad s(3)=0110
\]
into the even-weight subspace of \(\F_2^4\), and put \(r=1111\).  Let
\(\ell(y)\) be the first coordinate of \(s(y)\).  Define
\[
 h(e)=1_{D_0}(e)+\ell(\phi(e)),\qquad
 d(e)=s(\phi(e))+h(e)r.                            \tag{5}
\]
Both \(D_0\) and the support of the binary flow \(\ell\circ\phi\) are
Eulerian, so \(h\) is a binary flow.  Hence \(d\) is an
\(\F_2^4\)-flow.  For \(\phi(e)\ne0\), the two vectors
\(s(\phi(e))\) and \(s(\phi(e))+r\) are complementary weight-two
vectors.  Thus every \(d(e)\) has weight two.  Its first coordinate is
\[
 \ell(\phi(e))+h(e)=1_{D_0}(e).
\]
The four coordinate supports \(D_0,D_1,D_2,D_3\) of \(d\) therefore form
a four-cycle double cover of \(K\) with prescribed member \(D_0\).

Now take
\[
 A_1,\ A_2,\ D_1,\ D_2,\ D_3.
\]
An edge of \(M\) lies in \(A_1,A_2\) and nowhere else.  An edge of
\(D_0\) lies in exactly one of \(A_1,A_2\) and exactly one of
\(D_1,D_2,D_3\).  An edge outside \(D_0\cup M\) lies in neither
\(A_1,A_2\) and in exactly two of \(D_1,D_2,D_3\).  Every edge is covered
exactly twice.
\end{proof}

The conclusion and its matching/4-flow ingredients are consistent with,
and are a direct specialization of, the prior characterization in
\cite{hoffmann2013}.  The detailed proof is included only to make the
coordinate convention and the role of the two \(T\)-joins checkable.

\section{Circuit switches and the false domination statement}
\label{sec:switch}

Let \(a\ne0\), and let \(C\) be a circuit in \(G\).  Define
\[
 f^{C,a}(e)=
 \begin{cases}
   f(e)+a,&e\in C,\\
   f(e),&e\notin C.
 \end{cases}                                       \tag{6}
\]

\begin{lemma}[exact one-circuit exchange]\label{lem:switch}
If \(C\cap M_a=\varnothing\), then \(f^{C,a}\) is a nowhere-zero flow,
\(M'_a=M_a\), and, for \(b\notin\{0,a\}\),
\[
 M'_b=(M_b-C)\cup(M_{a+b}\cap C).                  \tag{7}
\]
\end{lemma}

\begin{proof}
At every vertex, the circuit has either zero or two incident edges.
Adding \(a\) on those two edges changes the conservation sum by
\(a+a=0\), so (6) remains a flow.  A switched edge becomes zero exactly
when it previously had value \(a\), excluded by the hypothesis.

No edge of \(M_a\) is switched.  No switched nonzero value can become
\(a\), since \(x+a=a\) would imply \(x=0\).  Thus \(M'_a=M_a\).  For
\(b\notin\{0,a\}\), an edge outside \(C\) retains value \(b\), while an
edge of \(C\) has new value \(b\) exactly when its old value was \(a+b\).
This is (7).
\end{proof}

The recent reconfiguration papers \cite{esperet2025,cranston2026} write
\(\mathcal F(G,A)\) for the graph whose vertices are nowhere-zero
\(A\)-flows and whose adjacent flows differ only on a circuit.  For
\(A=\F_2^3\), Lemma~\ref{lem:switch} describes exactly those adjacencies.
Indeed, if the difference of two \(A\)-flows is supported on a connected
circuit, conservation at a circuit vertex equates the difference values
on its two circuit edges; connectedness makes the difference constant
around the circuit.

Let \(\Good(G)\) be the set of good nowhere-zero \(\F_2^3\)-flows.

\medskip
\noindent
\textbf{One-circuit domination statement \((\mathsf D)\).}
\emph{For every finite connected bridgeless loopless cubic graph \(G\),
the set \(\Good(G)\) dominates \(\mathcal F(G,\F_2^3)\).}

Thus every nowhere-zero flow would either be good or have a good
neighbour obtained by one allowed connected-circuit switch.  This
statement is weaker than connectivity of the reconfiguration graph and
is not implied by the known absence of isolated flows.  More importantly,
Theorem~\ref{thm:maincountermodel} below proves that \((\mathsf D)\) is
false even for simple graphs.

\begin{proposition}[why \((\mathsf D)\) would have sufficed]
\label{prop:conditional}
If \((\mathsf D)\) were true, then every finite connected
bridgeless loopless cubic graph has a standard 5-CDC.
\end{proposition}

\begin{proof}
Every bridgeless graph has a nowhere-zero 6-flow by Seymour's theorem
\cite{seymour1981}, hence a nowhere-zero 8-flow.  Tutte's group-order
invariance for nowhere-zero flows \cite{tutte1954} supplies a nowhere-zero
\(\F_2^3\)-flow.  If it is good, apply
Proposition~\ref{prop:success}.  If it is bad,
\((\mathsf D)\) supplies a good adjacent flow in one move,
and the same proposition applies.
\end{proof}

Theorem~\ref{thm:maincountermodel} closes this particular route
negatively.  It says nothing negative about the conclusion of
Proposition~\ref{prop:conditional}; indeed, its countermodel graph has a
standard 5-CDC.

\section{A human-checkable bad supplied flow}
\label{sec:countermodel}

This section disproves the stronger assertion that every supplied
nowhere-zero \(\F_2^3\)-flow is already good.  The proof is finite and does
not depend on SAT or on the frontier program.

\begin{theorem}[countermodel to the every-flow claim]\label{thm:countermodel}
There is a connected simple bridgeless cubic graph \(G\) with a
nowhere-zero \(\F_2^3\)-flow \(f\) for which none of the seven nonzero
values is successful.  The same graph is 3-edge-colourable and has an
explicit standard 5-CDC.
\end{theorem}

\begin{proof}
Number the edges and assign values as follows.  The final column will be
used for a positive-control 3-edge-colouring.
\[
\begin{array}{c|c|c|c@{\qquad}c|c|c|c}
i&e_i&f(e_i)&c(e_i)&i&e_i&f(e_i)&c(e_i)\\ \hline
0&03&7&0& 8&28&3&2\\
1&06&6&1& 9&36&5&2\\
2&07&1&2&10&39&2&1\\
3&14&7&2&11&48&1&1\\
4&16&3&0&12&49&6&0\\
5&17&4&1&13&58&2&0\\
6&25&6&1&14&59&4&2\\
7&27&5&0&&&&
\end{array}                                                       \tag{8}
\]
Every vertex occurs in three rows, so \(G\) is cubic and simple.  The
vertex sequence
\[
 0,3,6,1,4,9,5,8,2,7,0
\]
is a Hamilton circuit.  Its edges lie on that circuit, and every remaining
edge is a chord whose union with either Hamilton path between its ends
contains a circuit.  Hence every edge lies on a circuit and \(G\) is
bridgeless.  Its graph6 encoding is \texttt{ICOef?kF?}.

At vertices \(0,\ldots,9\), the incident value triples are
\[
\begin{split}
 &(7,6,1),\ (7,3,4),\ (6,5,3),\ (7,5,2),\ (7,1,6),\\
 &(6,2,4),\ (6,3,5),\ (1,4,5),\ (3,1,2),\ (2,6,4).
\end{split}                                                       \tag{9}
\]
Each triple has xor zero and contains no zero, proving directly that
\(f\) is a nowhere-zero flow.  Its value classes are
\[
\begin{array}{c|c@{\qquad}c|c}
k&M_k&k&M_k\\ \hline
1&\{2,11\}&5&\{7,9\}\\
2&\{10,13\}&6&\{1,6,12\}\\
3&\{4,8\}&7&\{0,3\}\\
4&\{5,14\}&&
\end{array}                                                       \tag{10}
\]
and the endpoint pairs in every displayed class are disjoint.

It remains to prove that every class fails the criterion of
Lemma~\ref{lem:criterion}.  Fix \(k\), put a binary variable \(q_i\) on
edge \(i\), set \(q_i=0\) for \(i\in M_k\), and at every terminal set the
other two incident variables to one.  The only remaining conditions are
the ten vertex equations
\[
 \bigoplus_{i:e_i\ni v}q_i=0.                    \tag{11}
\]
Ordinary row reduction over \(\F_2\) gives the following complete reduced
systems.  The set \(U\) lists the variables not fixed before reduction.
\[
\begin{array}{c|l|l}
k&U&\text{reduced equations}\\ \hline
1&4,6,9,10,14&(q_4,q_6,q_9,q_{10},q_{14})=(0,0,1,0,1)\\
2&1,2,3,4,5,7&q_1+q_5=1,\ q_2+q_5=0,\ q_3=0,\
                  q_4+q_5=0,\ q_7=0\\
3&0,2,10,12,14&(q_0,q_2,q_{10},q_{12},q_{14})=(1,0,0,0,0)\\
4&0,1,8,9,11&0=1\\
5&3,11,12,13,14&0=1\\
6&5&0=1\\
7&6,7,8,13,14&0=1
\end{array}                                                       \tag{12}
\]
For an additional rank check, the coefficient/augmented ranks in the
last four rows are respectively \(4/5,4/5,1/2,4/5\).  Thus values
4 through 7 admit no forced-edge Eulerian set.  Values 1 and 3 have one
solution each, and value 2 has two solutions according to the free choice
of \(q_5\).  Their circuit components are exactly
\[
\begin{array}{c|l|c}
k&\text{circuit edge sets}&\text{terminal counts}\\ \hline
1&\{0,1,9\}\ ;\ \{3,5,7,8,12,13,14\}&1,3\\
2&\{0,1,9\}\ ;\ \{6,8,11,12,14\}&1,3\\
2&\{0,2,4,5,9\}\ ;\ \{6,8,11,12,14\}&1,3\\
3&\{0,1,9\}\ ;\ \{3,5,6,7,11,13\}&1,3
\end{array}                                                       \tag{13}
\]
Equations (11) make the completeness of (12) a direct five-variable
calculation in the largest case; substitution into (8) checks every
listed circuit.  Every possible circuit decomposition has odd terminal
counts on both components.  Lemma~\ref{lem:criterion} therefore excludes
two edge-disjoint \(T\)-joins for all seven \(M_k\).

Finally, the colour column of (8), in edge order, is
\[
 (0,1,2,2,0,1,1,0,2,2,1,1,0,0,2).               \tag{14}
\]
At every vertex the three incident colours are \(0,1,2\).  For each
colour, take the union of the other two colour classes.  These three
Eulerian edge sets cover every edge exactly twice; append two empty sets.
Thus \(G\) has a standard 5-CDC.
\end{proof}

The frozen artifact contains a second checker which, without using
Lemma~\ref{lem:criterion}, enumerates all \(T\)-joins in each complement.
It finds \(16,16,16,16,16,8,16\) joins for values \(1,\ldots,7\) and
minimum pairwise intersection one in every case.  This is a computational
cross-check of the human proof, not a replacement for it.

\section{The connected countermodel}
\label{sec:countermodel40}

We now refute \((\mathsf D)\).  The proof is a composition argument from
the ten-vertex block \(B\) of Theorem~\ref{thm:countermodel} and a second
ten-vertex graph with three visibly non-cyclable monochromatic edges.  No
SAT certificate is needed for the main construction.

Let \((G_1,f_1)\) and \((G_2,f_2)\) be vertex-disjoint cubic flow graphs,
and let
\[
 e_1=x_1y_1\in E(G_1),\qquad e_2=x_2y_2\in E(G_2)
\]
have the same value \(b\).  Their \emph{aligned two-sum} deletes
\(e_1,e_2\), adds \(x_1x_2,y_1y_2\), and gives both new edges value \(b\).
The other pairing of the four ends is equivalent for the arguments
below.  The two new edges are called the connectors.

\begin{lemma}[flow, structure, and circuit projection]
\label{lem:twosum-flow}
An aligned two-sum of connected bridgeless cubic flow graphs is again a
connected bridgeless cubic flow graph.  It is simple when the summands are
simple.  Every circuit of the sum uses zero or both connectors.  If it
uses both, replacing its path through either summand by that summand's
deleted edge gives a circuit in the graph before the sum.
If the summands have standard 5-CDCs, then so does the sum.
\end{lemma}

\begin{proof}
Every affected vertex loses and gains one incident edge of the same flow
value \(b\); hence conservation and degree three are preserved.  Deleting
a non-bridge from either connected summand leaves it connected, and the
two connectors join the two remaining graphs, so the sum is connected.
Each connector lies on a circuit formed from the two connectors and
paths between the deleted-edge ends in the two summands.  An old edge
lies on an old circuit; if that circuit used a deleted edge, replace the
deleted edge by the corresponding path through the other summand.  Thus
every edge of the sum lies on a circuit, proving bridgelessness.
Simplicity is immediate because the connectors join disjoint vertex
sets.

A circuit meets the displayed two-edge cut evenly, so it uses zero or
both connectors.  In the latter case, deleting the two connectors cuts
the circuit into one path in each summand, between the ends of the
deleted edge.  Restoring that edge closes either path to a circuit.  This
is the claimed projection.

For the cover assertion, use Proposition~\ref{prop:labels}.  Permute the
five coordinates of one summand so the two deleted edges have the same
two-subset label, and give both connectors that label.  At each affected
vertex one copy of that label is removed and one is added.  Exact-two
coverage and all five incidence parities are therefore preserved.
\end{proof}

\begin{lemma}[non-cyclability is inherited by port replacement]
\label{lem:twosum-ports}
Let \(P\subseteq E(H)\) contain at least two edges, and suppose no circuit
of \(H\) contains all of \(P\).  Replace every \(p\in P\) by an aligned
two-sum with a separate graph.  Every circuit of the resulting graph is
disjoint from one whole attached graph and its two connectors.
\end{lemma}

\begin{proof}
If the circuit lies wholly in one attached graph, any other attached
graph is missed.  Otherwise project every traversed attachment back to
its replaced edge using Lemma~\ref{lem:twosum-flow}.  The result is a
circuit of \(H\), so it omits some \(p\in P\).  The corresponding
attached graph and both connectors were not traversed.
\end{proof}

\begin{lemma}[one unchanged bad block blocks every value]
\label{lem:twosum-block}
Attach a copy of \(B\) by an aligned two-sum at an edge \(xy\) of value
\(b\).  Suppose a later modification leaves the flow on the \(B\)-side
and on both connectors unchanged.  Then no value \(k\ne0\) is successful
in the resulting global flow.
\end{lemma}

\begin{proof}
Suppose instead that \(J_1,J_2\) are edge-disjoint global
\(\bd M_k\)-joins in the complement of \(M_k\).

If \(k=b\), both connectors are removed with \(M_k\).  Restricting
\(J_1,J_2\) to the isolated \(B\)-side gives two edge-disjoint
\(\bd M_k(B)\)-joins in \(B-M_k(B)\), contradicting
Theorem~\ref{thm:countermodel}.

Let \(k\ne b\).  Both connectors remain.  For either \(J_i\), the parity
sum over all vertices on the \(B\)-side shows that it uses an even number
of connectors: the terminals in that shore are exactly
\(\bd M_k(B)\), a set of even cardinality.  Hence it uses zero or two
connectors.  If it uses none, its restriction is a
\(\bd M_k(B)\)-join.  If it uses both, delete the connectors and restore
the port edge \(xy\); the resulting edge set is again a
\(\bd M_k(B)\)-join.  Because \(J_1,J_2\) are edge-disjoint, at most one
of them uses both connectors, so at most one projection uses \(xy\).
The two projected joins are therefore edge-disjoint, again contradicting
the defining property of \(B\).
\end{proof}

\subsection{A ten-vertex non-cyclable base}

Let \(H\) have the following edge table.  The column \(c(e)\) is a proper
three-edge-colouring, and the flow value is \(1,2,3\) on colours
\(0,1,2\), respectively.
\[
\begin{array}{c|c|c|c@{\qquad}c|c|c|c}
i&e_i&c(e_i)&f(e_i)&i&e_i&c(e_i)&f(e_i)\\ \hline
0&01&0&1& 8&67&2&3\\
1&02&1&2& 9&18&2&3\\
2&03&2&3&10&48&1&2\\
3&23&0&1&11&58&0&1\\
4&45&2&3&12&29&2&3\\
5&46&0&1&13&39&1&2\\
6&56&1&2&14&79&0&1\\
7&17&1&2&&&&
\end{array}                                                       \tag{15}
\]
Its graph6 record is \texttt{It?GYDKKO}.  The table gives every vertex
one edge of each colour, so \(H\) is cubic and the three incident flow
values xor to zero.  It is connected: from \(0\) one reaches
\(1,2,3\), then \(7,8,9,6\), and then \(4,5\).  Every edge lies in the
two-colour 2-factor obtained by omitting the third colour, so \(H\) is
bridgeless.

Put
\[
 P=\{1,6,7\}=\{02,56,17\}.                                      \tag{16}
\]
All three edges have value \(2\), but no connected circuit contains all
three.  Indeed, suppose a circuit \(C\) contains \(P\).  It has edges
outside the triangle on \(\{0,2,3\}\), so it crosses that triangle's
three-edge boundary in two edges.  If \(01\notin C\), then \(29,39\in C\);
together with \(02\), the forced degrees at \(0,2,3,9\) make
\(0\,2\,9\,3\,0\) the whole circuit, a contradiction.  Hence
\(01\in C\).

Apply the same argument to the triangle on \(\{4,5,6\}\).  If
\(67\notin C\), the forced edges make \(8\,4\,6\,5\,8\) the whole
circuit.  Thus \(67\in C\), and exactly one of \(48,58\) lies in \(C\).
At vertex \(7\), the edges \(17,67\) exclude \(79\).  At vertex \(1\),
the edges \(17,01\) exclude \(18\).  Vertex \(8\) is consequently
incident with exactly one edge of \(C\), namely one of \(48,58\), which
is impossible for a circuit.  This proves the non-cyclability of \(P\)
without computation.

\begin{theorem}[40-vertex countermodel to one-circuit domination]
\label{thm:maincountermodel}
There is an explicit simple connected bridgeless cubic graph \(G\) on
40 vertices and 60 edges with a nowhere-zero \(\F_2^3\)-flow \(f\) such
that \(f\notin\Good(G)\) and every allowed switch of \(f\) on one
connected circuit also lies outside \(\Good(G)\).  Moreover, \(G\) has
an explicit three-cycle double cover.  Thus \((\mathsf D)\) is false,
but \(G\) is not a counterexample to the five-cycle double cover
conjecture.
\end{theorem}

\begin{proof}
For each of the three edges \(p\in P\), replace \(p\) by an aligned
two-sum with a fresh copy of \(B\), using edge \(10\) of \(B\) as the
port.  Both edges have value \(2\).  Repeated use of
Lemma~\ref{lem:twosum-flow} proves that the resulting \(G,f\) is simple,
connected, bridgeless, cubic, and nowhere zero.  The counts are
\[
 |V(G)|=10+3\cdot10=40,\qquad |E(G)|=15+3\cdot15=60.
\]

By Lemma~\ref{lem:twosum-ports}, every connected circuit \(C\) of \(G\)
misses one whole copy of \(B\) and its connectors.  An allowed switch on
\(C\) leaves that block unchanged, so Lemma~\ref{lem:twosum-block}
shows that all seven values of the switched flow remain unsuccessful.
The original flow is bad by the same lemma, using any attached block.

For the positive assertion, edge \(10\) of \(B\) and all three edges of
\(P\) have colour \(1\) in the displayed proper three-edge-colourings.
Give both connectors colour \(1\).  At each affected vertex an edge of
colour \(1\) is replaced by another, so the final graph remains properly
three-edge-coloured.  For each colour, take the union of the other two
colour classes.  These three Eulerian edge sets form a cycle double
cover, each with 40 edges.  Appending two empty sets gives a standard
5-CDC.
\end{proof}

\medskip
\noindent
\textbf{Graph metadata.}
The nauty-canonical graph6 encoding is obtained by concatenating the
following two displayed lines:
\begin{center}
\small
\texttt{\detokenize{gt?GG?@?GB?@?D??@O?K???_?CGA_??A???@??OO??C??@_??C???G?_???G???D???C??}}\\
\texttt{\detokenize{??_G???G????OC????B????K?????G????@@O????KO????BA?GC??@O??O??}}.
\end{center}
NetworkX and nauty \texttt{planarg} report that \(G\) is planar, and its
vertex- and edge-connectivity are both two.  These are metadata about the
route countermodel, not evidence against five-CDC; the explicit 3-CDC
above proves the opposite for this graph.

The construction has a visible nontrivial two-edge cut at each block
attachment, so it does not address \((\mathsf D)\) restricted to
cyclically 4-edge-connected graphs.

The complete construction, flow, three-cover labels, checker, and frozen
replay are in
\begin{center}
\texttt{search/connected-one-switch-countermodel-40v-20260726/}.
\end{center}
The checker does not import the producer.  It reconstructs the three
two-sums; verifies simplicity, cubicity, connectedness, every single-edge
deletion, and every flow equation; enumerates all 30 connected circuits
of \(H\) and all 6,780 connected circuits of \(G\); confirms that every
expanded circuit misses an entire bad block; re-enumerates all local
\(T\)-joins; and checks the three-cover coordinate sizes
\[
 40,\quad40,\quad40,\quad0,\quad0.
\]
This is an AI-assisted mechanical replay, not independent human review.

\section{A potential-compatible pure-merge countermodel}
\label{sec:puremerge}

We now study a different, narrower route from a fixed
\(\F_2^3\)-flow to five Eulerian coordinates.  The potential construction
below is the flow-to-pair-label construction in the July 2026
cycle-double-cover proof, as presented by Oum and Kintali
\cite[Section~5]{oum2026} and \cite{kintali2026}.  We use
\emph{potential-compatible} only as local shorthand for covers arising
from these equations; it is not terminology attributed to those authors.

\subsection{Potentials, pair labels, and pure merges}

Let \(G\) be cubic and let
\(\phi:E(G)\to\F_2^3\setminus\{0\}\) be a flow.  A family
\((t_v:v\in V(G))\) of vertex potentials is \emph{compatible} when, for
every edge \(e=uv\),
\[
 t_u+t_v\in d_e+\langle\phi(e)\rangle,\qquad
 d_e=\phi(f_u)+\phi(f_v),                                      \tag{P}
\]
where \(f_u\ne e\) and \(f_v\ne e\) are incident with \(u\) and \(v\).
The choice of either other edge at an endpoint is immaterial because the
three incident flow values sum to zero.  A compatible potential defines
the two-element label
\[
 P_e=t_u+\phi(f_u)+\langle\phi(e)\rangle
     =t_v+\phi(f_v)+\langle\phi(e)\rangle.                       \tag{L}
\]
For \(s\in\F_2^3\), let \(A_s=\{e:s\in P_e\}\).  If \(e,f,g\) are the
three edges at \(v\), then their labels at that end are
\[
 \{t_v+\phi(f),t_v+\phi(g)\},\quad
 \{t_v+\phi(e),t_v+\phi(g)\},\quad
 \{t_v+\phi(e),t_v+\phi(f)\}.
\]
Every coordinate consequently occurs zero or twice at \(v\).  Thus all
eight \(A_s\) are Eulerian and cover each edge exactly twice.  We call
such a family a \emph{potential-compatible eight-cover for the fixed flow
\(\phi\)}.

A \emph{pure merge} assigns the eight coordinates to at most five
classes and takes the symmetric difference of the \(A_s\) inside each
class.  It changes no edge within a coordinate before this final merging
step.  The \emph{coordinate co-occurrence graph} \(Q(P)\) has vertex set
\(\F_2^3\), with \(ab\) an edge when \(P_e=\{a,b\}\) for some
\(e\in E(G)\).

\begin{proposition}[pure-merge colouring equivalence]
\label{prop:pure-colour}
A potential-compatible eight-cover admits a pure merge into at most five
Eulerian edge-subsets, still covering every edge exactly twice, if and
only if \(Q(P)\) is properly five-colourable.
\end{proposition}

\begin{proof}
Colour each coordinate by its merge class.  An edge with label
\(\{a,b\}\) survives in two output coordinates exactly when \(a\) and
\(b\) have different colours; if their colours agree, its two
occurrences cancel in the same symmetric difference.  Hence a valid
merge gives a proper colouring of \(Q(P)\).  Conversely, use the colour
classes of a proper colouring as merge classes.  Symmetric differences
of Eulerian edge sets are Eulerian, and every original edge survives
once in each of two different classes.
\end{proof}

This equivalence concerns only a pure identification of the eight
coordinates.  It says nothing against a five-cycle double cover obtained
by a different construction.

\subsection{A potential-rigid twelve-vertex cap}

Let \(R\) be the cubic graph with graph6 record
\(\texttt{K??FEaKR@oE\_}\).  The following table freezes its edge order,
flow, proper edge-colouring \(c\), the vector \(d_e\) from (P), and one
compatible pair label from (L).
\[
\begin{array}{c|c|c|c|c|c@{\qquad}c|c|c|c|c|c}
i&e_i&\phi(e_i)&c(e_i)&d_i&P_i&
i&e_i&\phi(e_i)&c(e_i)&d_i&P_i\\ \hline
0&06&1&0&6&23& 9&37&5&2&5&36\\
1&07&2&1&6&13&10&3,10&7&1&6&34\\
2&08&3&2&4&12&11&3,11&2&0&3&46\\
3&16&4&1&6&26&12&48&5&0&2&14\\
4&17&7&0&6&16&13&49&1&1&6&01\\
5&19&3&2&5&12&14&4,10&4&2&6&04\\
6&26&5&2&2&36&15&58&6&1&1&24\\
7&2,10&3&0&2&03&16&59&2&0&5&02\\
8&2,11&6&1&7&06&17&5,11&4&2&0&04
\end{array}                                                       \tag{18}
\]
Here pairs such as \(23\) mean \(\{2,3\}\), not an integer.
The three incident flow values xor to zero at every vertex, and the
colour column is proper.  In particular, every edge of \(R\) lies in a
two-colour 2-factor, so \(R\) is bridgeless; connectedness is immediate
from the displayed incidence table.

Use edge \(1=07\) as the port and delete it.  We next give a direct
row-reduction proof that the retained compatibility equations determine
all twelve potentials up to adding one common vector.

Fix \(t_0=0\).  On the spanning tree whose edges, in order, are
\[
 0,\ 2,\ 3,\ 6,\ 12,\ 15,\ 4,\ 5,\ 7,\ 8,\ 9,
\]
write \(t_u+t_v=d_e+\lambda_j\phi(e)\) for the \(j\)-th displayed tree
edge \(e=uv\).  These tree equations express all \(t_v\) in the eleven
bits \(\lambda_0,\ldots,\lambda_{10}\).  Substitution in the six
remaining internal edges \(10,11,13,14,16,17\), followed by elementary
xor row operations, gives
\[
\begin{array}{rcl@{\qquad}rcl}
\lambda_3+\lambda_{10}&=&0&
\lambda_2+\lambda_6+\lambda_9&=&0\\
\lambda_2+\lambda_8&=&1&
\lambda_1+\lambda_7&=&0\\
\lambda_6&=&1&
\lambda_2+\lambda_5&=&1\\
\lambda_2+\lambda_4&=&1&
\lambda_0+\lambda_1+\lambda_3&=&0\\
\lambda_1+\lambda_2&=&0&
\lambda_1&=&1\\
\lambda_0&=&1.&&&
\end{array}
\]
The unique solution and its back-substitution are
\[
\begin{split}
(\lambda_0,\ldots,\lambda_{10})
  &=(1,1,1,0,0,0,1,1,0,0,0),\\
(t_0,\ldots,t_{11})
  &=(0,5,5,1,5,6,7,4,7,3,7,2).
\end{split}                                                       \tag{19}
\]
Thus the 37 scalar equations---two for each retained edge and three
fixing \(t_0\)---have rank 36.

Computing (L) from (19), the internal labels and the dangling port label
\(13\), calculated from either end of the deleted port, use exactly
\[
\begin{gathered}
01,\ 02,\ 03,\ 04,\ 06,\qquad
12,\ 13,\ 14,\ 16,\\
23,\ 24,\ 26,\qquad
34,\ 36,\qquad
46.
\end{gathered}                                                    \tag{20}
\]
These are all 15 pairs on \(\{0,1,2,3,4,6\}\), so the co-occurrence
graph contains \(K_6\).

\begin{lemma}[rigid-cap obstruction]\label{lem:rigid-cap}
Suppose a copy of \(R-e_1\), with the displayed internal flow and two
dangling flow-\(2\) connectors, occurs in a larger cubic flow graph.
Every compatible global potential forces a coordinate \(K_6\) on the
cap and its connectors.  Hence no compatible eight-cover for that fixed
global flow admits a pure merge into at most five coordinates.
\end{lemma}

\begin{proof}
Restrict a compatible global potential to the cap.  The retained
equations just solved show that it is (19) plus one common vector
\(a\in\F_2^3\).  Formula (L) at the cap ends gives the translated dangling
label \(13+a\) on the connectors.  Thus (20) is translated by \(a\), and
still induces a \(K_6\).  Proposition~\ref{prop:pure-colour} completes
the proof.
\end{proof}

\subsection{The 46-vertex one-switch countermodel}

\begin{theorem}[46-vertex pure-merge countermodel]
\label{thm:puremerge}
There is an explicit simple connected bridgeless cubic graph \(G^\star\)
on 46 vertices and 69 edges, with a nowhere-zero
\(\F_2^3\)-flow \(\phi\), such that:
\begin{enumerate}
  \item no potential-compatible eight-cover for \(\phi\) admits a pure merge
        into at most five Eulerian edge-subsets; and
  \item for every connected circuit \(C\) and every nonzero
        \(a\in\F_2^3\) for which
        \(\phi'=\phi+a\chi_C\) remains nowhere zero, no
        potential-compatible eight-cover for \(\phi'\) admits such a pure
        merge.
\end{enumerate}
Moreover, \(G^\star\) is Tait-colourable and has an explicit
three-cycle double cover.  It is therefore not a counterexample to the
five-cycle double cover conjecture.
\end{theorem}

\begin{proof}
Start from the base \(H\) in (15).  For each
\(p\in P=\{02,56,17\}\), replace \(p\) by an aligned two-sum with a
fresh copy of \(R\), using edge \(1=07\) as the cap port.  Both the base
ports and the cap port have flow value \(2\).  The resulting counts are
\[
 |V(G^\star)|=10+3\cdot12=46,\qquad
 |E(G^\star)|=15+3\cdot18=69.
\]
Lemma~\ref{lem:twosum-flow}, together with the incidence table above,
proves that the displayed graph and flow are simple, connected,
bridgeless, cubic, and nowhere zero.

The original flow has a compatible potential: put \(t_v=0\) on the base
and use (19) on each cap.  The base pair labels are \(23,13,12\) on
colours \(0,1,2\), respectively, so both sides assign \(13\) to every
connector.  Lemma~\ref{lem:rigid-cap} applied to any cap proves (i).

Let \(C\) be a connected circuit of \(G^\star\).  The non-cyclability
proof following (16), equivalently Lemma~\ref{lem:twosum-ports}, shows
that \(C\) leaves one whole cap and both its connectors untouched.
For every allowed switch \(\phi'=\phi+a\chi_C\), the flow on that cap
and its connectors is unchanged.  Restricting any potential compatible
with \(\phi'\) to the untouched cap and applying
Lemma~\ref{lem:rigid-cap} gives a forced \(K_6\).  This proves (ii).

Finally, the base ports and cap port all have colour \(1\) in the proper
edge-colourings in (15) and (18).  Give the connectors colour \(1\).
This produces a proper three-edge-colouring of \(G^\star\).  If its
colour classes are \(M_0,M_1,M_2\), then
\[
 E(G^\star)\setminus M_0,\quad
 E(G^\star)\setminus M_1,\quad
 E(G^\star)\setminus M_2
\]
are Eulerian and cover every edge exactly twice.  Each has 46 edges;
append two empty sets for a standard 5-CDC with sizes
\(46,46,46,0,0\).
\end{proof}

The construction has three visible nontrivial two-edge cuts, one at each
cap.  It therefore does not refute a pure-merge statement restricted to
cyclically 4-edge-connected cubic graphs.  The frozen construction
metadata reports that \(G^\star\) is nonplanar; this fact is not used in
the proof, and the independent checker does not recheck planarity.

\medskip
\noindent
\textbf{Exact encoding and check.}
The nauty-canonical graph6 encoding is obtained by concatenating:
\begin{center}
\small
\texttt{\detokenize{ms????@?G@???P?D?@_?o?_?A?????O?CC?I@?@_O?W??Q???W??_?O????G???K????G???}}\\
\texttt{\detokenize{?OA????W????????E????I?????A_????S????B????????????????????C?????WO????@}}\\
\texttt{\detokenize{Q?????@`?????@__????B?O????B??}}.
\end{center}
The complete construction and SAT-free proof package are in
\begin{center}
\texttt{search/fano-pure-merge-one-switch-countermodel-46v-20260726/}.
\end{center}
Its independently written standard-library checker reconstructs the
graph from the ten- and twelve-vertex constants; checks all 69
single-edge deletions, every flow and parity equation, and the explicit
three-cover; independently row-reduces the cap system to rank 36 and
checks the forced \(K_6\); and enumerates all 30 base circuits.  As a
further cross-check, it enumerates all 128 globally compatible gauged
potentials for the original flow: every co-occurrence graph contains a
\(K_6\), and none is five-colourable.  The checker and this audit remain
AI-assisted rather than independent human verification.

\section{Finite reconfiguration frontier}
\label{sec:computation}

These searches preceded the construction of
Theorem~\ref{thm:maincountermodel}.  They are retained because they show
how far the false statement survives finite testing, not because they
support it as a universal claim.  We state only what the artifacts
certify.  They test the successful-value condition and do not test the
distinct pure-merge property in Theorem~\ref{thm:puremerge}.  A \emph{hard
graph} in this section is connected, simple, bridgeless, cubic, and not
3-edge-colourable.  A \emph{strict snark} is additionally cyclically
4-edge-connected and has girth at least five.  The flow enumeration quotients by
\(\mathrm{GL}(3,2)\).  A bad flow uses all seven nonzero values, since an
empty value class has \(M_k=\varnothing\) and is immediately successful.
The retained normalization first fixes the ordered value line at a root
vertex and then removes its residual stabilizer.

\medskip
\noindent
\textbf{Computational Observation.}
\emph{The retained programs report the following exact finite totals.}
\[
\begin{array}{lrrrr}
\toprule
\text{scope}&\text{flow orbits}&\text{bad}&
\text{reported repair}&\text{unrepaired}\\
\midrule
\text{all bridgeless cubic, orders }4\text{--}10
 &3{,}295&5&5&0\\
\text{all bridgeless cubic, order }12
 &73{,}152&304&304&0\\
\text{all bridgeless cubic, order }14
 &2{,}216{,}590&18{,}152&18{,}152&0\\
\text{26 hard graphs, order }16
 &600{,}440&10{,}413&10{,}413&0\\
\text{179 hard graphs, order }18
 &21{,}324{,}510&757{,}221&757{,}221&0\\
\text{1,388 hard graphs, order }20
 &847{,}539{,}220&49{,}752{,}091&49{,}752{,}091&0\\
\text{20 strict snarks, order }22
 &63{,}251{,}330&279{,}836&279{,}836&0\\
\text{38 retained strict snarks, order }24
 &617{,}079{,}260&5{,}161{,}169&5{,}161{,}169&0\\
\text{one retained girth-6 snark, order }28
 &432{,}613{,}280&3{,}400{,}482&3{,}400{,}482&0\\
\bottomrule
\end{array}                                                       \tag{17}
\]
Here the first five retained producers test arbitrary nonempty binary
cycles, which may have more than one circuit component.  Such a switch
decomposes into a valid sequence of circuit switches because the switched
cycle avoids the \(a\)-class, which remains fixed throughout.  But a
one-binary-cycle repair is not by itself a certificate of
\emph{one-circuit adjacency} from the starting flow.

The complete order-20 hard-graph computation was repeated with the switch
restricted to a connected circuit.  The elementary-only C++ ledger again
contains all \(847{,}539{,}220\) orbit representatives and reports all
\(49{,}752{,}091\) bad representatives repaired in one circuit, with zero
unrepaired.  The same elementary-circuit enumerator covers the complete
frozen sets of 20 strict order-22 snarks and 38 strict order-24 snarks.
At order 22, among \(63{,}251{,}330\) flow orbits, all \(279{,}836\)
initially bad representatives are repaired.  At order 24, among
\(617{,}079{,}260\) flow orbits, \(611{,}918{,}091\) are initially good
and all \(5{,}161{,}169\) initially bad representatives are repaired.
Neither scope contains a switch-local bad flow.  The order-24 source is
the complete output of one retained Snarkhunter 2.0b command and is
frozen with the ledger.  The verifier checks all 38 distinct graph
identities and their strict-snark premises, but does not independently
regenerate the canonical list; census completeness therefore relies on
Snarkhunter.  These last three rows are direct finite tests of the
now-refuted statement \((\mathsf D)\), with no extrapolation beyond the
displayed scopes.  No analogous full elementary-only ledger is claimed
here for the earlier rows.  The
\(40\)-vertex countermodel demonstrates that these zero-obstruction
totals are a finite phenomenon.

Separately, the enumerator exhausts \(432{,}613{,}280\) flow orbits on
one retained cyclically 4-edge-connected girth-6 snark of order 28.  All
\(3{,}400{,}482\) initially bad representatives are repaired in one
circuit.  This is one graph, not a complete order-28 graph census.

\subsection*{Artifacts and checks}

The ten-vertex block package is
\begin{center}
\texttt{search/fano-value-class-flow-countermodel-20260726/},
\end{center}
the connected countermodel package is
\begin{center}
\texttt{search/connected-one-switch-countermodel-40v-20260726/},
\end{center}
the pure-merge countermodel package is
\begin{center}
\texttt{search/fano-pure-merge-one-switch-countermodel-46v-20260726/},
\end{center}
and the finite-frontier package is
\begin{center}
\texttt{search/fano-flow-one-switch-frontier-20260726/}.
\end{center}
From the repository root, the principal checks are
\begin{verbatim}
(cd search/fano-value-class-flow-countermodel-20260726 &&
 python3 independent_checker.py &&
 shasum -a 256 -c SHA256SUMS)

(cd search/fano-flow-one-switch-frontier-20260726 &&
 python3 verify.py &&
 shasum -a 256 -c SHA256SUMS)

(cd search/connected-one-switch-countermodel-40v-20260726 &&
 python3 independent_checker.py &&
 shasum -a 256 -c SHA256SUMS)

(cd search/fano-pure-merge-one-switch-countermodel-46v-20260726 &&
 python3 independent_checker.py)
\end{verbatim}
On the retained environment the first three package manifests pass, and
the fourth command reproduces the frozen 46-vertex checker output.  The block
checker reconstructs graph6 identity, cubicity, connectedness,
bridgelessness by every single-edge deletion, flow conservation, all seven
matching supports, literal \(T\)-join enumeration, the colouring, and the
resulting five-cover.

For the frontier, a separately written C++ implementation matches the
complete Python order-14 totals.  It is also the producer for the two
order-20 ledgers and the connected-circuit order-22, order-24, and
one-graph order-28 ledgers.  The Python
verifier checks the frozen aggregate identities and semantically
reconstructs the first retained bad-flow repair on 432 graph rows.  It
does not re-enumerate every order-16, order-18, order-20, order-22, or
order-24 flow orbit, nor the retained order-28 graph.

\subsection*{Limitations}

\begin{enumerate}
  \item The finite frontier does not contradict
        Theorem~\ref{thm:maincountermodel}; its tested graph orders and
        scopes simply do not include the constructed obstruction.  It
        concerns a different property from
        Theorem~\ref{thm:puremerge}.
  \item At orders 16, 18, and 20, the graph scope contains only hard
        graphs, and at orders 22 and 24 only strict snarks, not every
        bridgeless cubic graph.  This matters because
        Theorems~\ref{thm:maincountermodel} and \ref{thm:puremerge} are
        both 3-edge-colourable.
        The order-24 census additionally relies on the retained
        Snarkhunter run rather than an independent second generation.
        The order-28 row is explicitly only one graph.
  \item The programs and their cross-checks were all AI-assisted.  They
        are implementation diversity, not independent human validation.
  \item The full positive model or repair witness for every one of the
        hundreds of millions of orbit decisions is not retained as a
        separately checkable certificate archive.
  \item The computation uses simple graphs.  The structural statements
        above allow loopless cubic multigraphs, but no finite conclusion
        is extrapolated to untested parallel-edge instances.
  \item No orientability condition is encoded or inferred.
\end{enumerate}

\section{What remains open}

Theorem~\ref{thm:maincountermodel} ends the proposed one-connected-circuit
domination route, and Theorem~\ref{thm:puremerge} ends the proposed route
of applying one circuit switch and then purely merging the eight
coordinates supplied by potential-compatible labels.  Neither theorem
produces a graph without a 5-CDC: both graphs have displayed 3-CDCs.
The general five-cycle double cover conjecture therefore remains exactly
as open as before this work.  In particular, the second theorem does not
exclude changing coordinates by operations more general than a pure
merge, choosing a different flow after several switches, or constructing
a 5-CDC without the potential eight-cover.

Both constructions have nontrivial two-edge cuts.  Hence neither theorem
settles the corresponding one-circuit assertion for cyclically
4-edge-connected cubic graphs.  No reduction from either restricted
assertion to the general 5-CDC conjecture is claimed here.

The structural statements in Sections~\ref{sec:labels} and
\ref{sec:value} may still be useful independently.  A successful value
is a sufficient condition for a cover, and the four canonical
\(\bd M_k\)-joins retain their tetrahedral intersection pattern.  What
fails is the claim that one circuit switch must reach such a value.  Any
future proof strategy must either use longer/global reconfiguration,
combine information from unsuccessful values without requiring one to
become successful, or leave this flow framework.  For the pure-merge
route, a future argument would have to permit a richer transformation
than merely identifying coordinates, use more than one local flow
switch, or establish a restricted theorem on a domain not covered by the
cap construction.

\section*{AI-use statement}

The page-one disclosure is part of the mathematical record of this draft
and should remain in every derivative version.  OpenAI Codex agents
materially contributed to the definitions selected for study, conjecture,
proof search, the finite witnesses, both two-sum compositions, the
potential-rigid cap, source code, SAT encoding, checkers, experiments,
literature search, and exposition.
A second Codex agent performed an adversarial audit, but that is not
independent human peer review.  The recent flow-reconfiguration
framework is attributed to \cite{esperet2025,cranston2026}; the
matching/4-flow criterion is attributed to \cite{hoffmann2013}; and the
potential construction is attributed to \cite{oum2026,kintali2026}.  The two-sum
proofs, the three-edge non-cyclability argument, and the rigid-cap row
reduction are written for direct human checking, but no independent
human check is asserted to have occurred.  Any human authors should
record their verification, independent reconstruction, and rerun history
before considering public submission.

\bibliographystyle{plain}
\bibliography{references}

\end{document}
